\section{Описание методов}

\subsection{Алгоритмическая модель}

Первый стенд сравнивает BFS, алгоритм Дейкстры, двунаправленный алгоритм Дейкстры, A*, алгоритмы Беллмана--Форда и Флойда--Уоршелла. Набор включает 12 типов топологий: линию, кольцо, звезду, решётку, случайный граф, плотную mesh-сеть, две области с узким соединением, безмасштабную сеть, малый мир, дерево, fat-tree и dragonfly. Последние топологии используются для исследования путевой избыточности и не интерпретируются как геометрические модели размещения LoRa-устройств.

Для каждой топологии подготовлено 30 графов, всего 360. Размер зависит от семейства: для линии, кольца и звезды генератор задаёт от 20 до 300 вершин, для решётки --- от 4 до 20 строк и столбцов, для двоичного дерева --- высоту от 3 до 7. Исходные веса --- целые числа от 1 до 50. Все алгоритмы получают один и тот же граф и пару вершин. Для взвешенных методов проверяется совпадение стоимости пути с опорным решением. У A* корректность зависит от допустимости эвристики; отсутствие информативной геометрической оценки ограничивает возможное ускорение.

Время выполнения характеризует конкретную программную реализацию и вычислительную машину. Число посещённых вершин и обработанных рёбер даёт дополнительную характеристику объёма работы. Измерения на рабочей станции не являются прямым измерением задержки или расхода энергии микроконтроллера. Алгоритм Флойда--Уоршелла решает более широкую задачу всех пар, что учитывается при интерпретации его затрат на один запрос.

\subsection{Дискретно-событийная модель радиоканала}

Второй стенд моделирует создание сообщений, ожидание доступа к среде, передачу и приём кадров, ретрансляцию и изменение состояния узлов. События обрабатываются в порядке модельного времени. Общая реализация физического уровня и доступа к среде используется для всех сравниваемых стратегий.

Длительность кадра $T_{\mathrm{air}}$ вычисляется по формуле Semtech~\cite{semtech}. Длительность символа равна
\begin{equation}
  T_{\mathrm{sym}}=\frac{2^{SF}}{BW}.
\end{equation}
Полное время включает преамбулу, заголовок, полезную нагрузку, проверочную сумму и избыточность кодирования. Поэтому при одинаковом числе байтов увеличение SF увеличивает занятость канала.

Потери распространения задаются моделью
\begin{equation}
\begin{split}
  PL(d)&=PL(d_0)+10n\log_{10}(d/d_0)+X,\\
  X&\sim\mathcal{N}(0,\sigma^2).
\end{split}
\end{equation}
В базовой конфигурации используются $PL(d_0)=127{,}41$~дБ при $d_0=40$~м, $n=2{,}08$ и $\sigma=3{,}57$~дБ, принятые по~\cite{bor}. Это параметры выбранной среды, а не универсальная характеристика городского канала.

По мощности передатчика, усилению антенн и потерям определяется RSSI, затем отношение сигнал/шум SNR. Чувствительность и порог SNR зависят от SF. Вблизи порога декодирование описывается вероятностной зависимостью. Для совпадающих во времени передач проверяется возможность захвата более сильного сигнала с порогом 6~дБ. Такой порог является упрощением модели интерференции.

Доступ к среде включает обнаружение занятости, случайную отсрочку и ограничение числа попыток. Рабочий цикл в базовом сценарии равен 1~\% и учитывается в скользящем окне 3600~с. Кадр, превышающий бюджет узла, отбрасывается. Значение 1~\% задаёт исследуемый сценарий; требования конкретного частотного поддиапазона и региона должны устанавливаться отдельно.

Для одноадресных кадров моделируются подтверждения и до трёх повторных передач. Широковещательные кадры не подтверждаются. Энергия рассчитывается по состояниям передачи, приёма, ожидания и сна:
\begin{equation}
  E_i=U\sum_k I_k t_{ik}.
\end{equation}
Используются постоянные токи состояний. Модель не воспроизводит нелинейный разряд батареи и индивидуальные характеристики всех аппаратных платформ.

\begin{table}[t]
\caption{Базовые параметры сетевого стенда}
\label{tab:params}
\centering
\small
\begin{tabular}{@{}p{0.58\columnwidth}r@{}}
\toprule
Параметр & Значение \\
\midrule
Число узлов & 25 \\
Коэффициент расширения & SF9 \\
Полоса & 125 кГц \\
Скорость кодирования & 4/5 \\
Мощность передатчика & 14 дБм \\
Полезная нагрузка / заголовок & 32 / 16 байт \\
Интенсивность для всей сети & 6 сообщений/мин \\
Основное модельное время & 1200 с \\
Завершение оставшихся передач & 120 с \\
Рабочий цикл & 1\,\% \\
Ёмкость батареи & 2000 мА$\cdot$ч \\
\bottomrule
\end{tabular}
\end{table}

\subsection{Сравниваемые стратегии}

\textit{AODV.} Используется модель обнаружения пути сообщениями запроса и ответа. Она позволяет исследовать стоимость поиска и восстановления маршрута. Это упрощённая реализация механизмов AODV, а не проверка полного соответствия RFC~3561.

\textit{OLSR.} Моделируются периодические сообщения о соседях и топологии и распространение с использованием MPR. Цель --- оценить стоимость поддержания маршрутов при дефиците эфира. Полнота промышленной реализации RFC~3626 не предполагается.

\textit{ManagedFlood.} Узел ретранслирует ранее не встречавшееся сообщение с ограничением числа переходов. Время ожидания зависит от SNR: меньший запас сигнала даёт более раннюю попытку. Если до неё принят дубликат, запланированная передача отменяется. В базовой реализации предел равен восьми переходам.

\textit{EpidemicGossip.} Моделируется вероятностное лавинное распространение: после первого приёма сообщение передаётся с вероятностью 0,8, предел --- десять переходов. В отличие от ManagedFlood подавление ожидающей передачи по дубликату отсутствует. Здесь термин обозначает gossip-схему и не подразумевает полноценный DTN-протокол с обменом буферами при встречах узлов.

\textit{Ideal.} Эталон мгновенно получает доступный граф и выбирает путь минимальной стоимости без служебного обнаружения. Передачи по выбранному маршруту подчиняются тем же ограничениям канала. Эталон относится к одноадресной передаче по одному пути и не является верхней границей PDR для произвольного способа распространения.

\subsection{Связь уровней через физические веса}

Алгоритмический стенд повторно применяется к графам, построенным сетевой моделью. Для ребра задаётся стоимость
\begin{equation}
  w_{uv}=\frac{T_{\mathrm{air}}}{p_{uv}},
  \label{eq:weight}
\end{equation}
где $p_{uv}$ --- вероятность приёма кадра при среднем состоянии линии без текущих коллизий. Стоимость имеет размерность времени и характеризует ожидаемые затраты на успешную передачу в упрощённой модели независимых попыток. При ограниченном числе повторений, наличии подтверждений и конкуренции за эфир она служит маршрутной оценкой, а не точным временем доставки.

В реализации различаются граф физически слышимых соседей и граф каналов для маршрутизации. Во второй включаются линии с $p_{uv}\geq0{,}7$; по нему строит путь Ideal. Порог и множество допустимых каналов влияют на сравнение: лавинное распространение может использовать приём вне заранее выбранного одноадресного пути.

Для проверки переноса весов используются три типа размещения, три значения seed и до 12 пар вершин на граф. На этих графах повторяются алгоритмы первого стенда и проверяется стоимость решения~\eqref{eq:path}. Эта проверка связывает уровни задачи, но сама по себе не доказывает превосходство какой-либо стратегии передачи.

\subsection{Сценарии и статистическая обработка}

Основное сравнение включает статическую сеть, случайные отказы с восстановлением и одновременное отключение 35~\% узлов. В потоке сообщений интервалы поступления имеют экспоненциальное распределение. Пары отправителя и получателя выбираются среди связанных в исходном маршрутном графе узлов. Поэтому PDR характеризует передачу для этого набора пар, а не покрытие всех возможных пар на территории.

Для основного сравнения задаются десять seed от 42 до 51. При одинаковом seed все стратегии получают одинаковое размещение, поток прикладных сообщений и расписание внешних отказов. Радиособытия и потребление случайных чисел после начала работы могут различаться из-за поведения протоколов. Для каждого показателя рассчитываются среднее и приближённый 95-процентный доверительный интервал $\bar{x}\pm1{,}96s/\sqrt{r}$, где $s$ --- выборочное стандартное отклонение, а $r$ --- число независимых прогонов. При малом числе повторений точность нормального приближения ограничена. Наблюдения внутри одной сети не считаются независимыми повторениями эксперимента.

Дополнительные серии изменяют число узлов от 15 до 80, SF от 7 до 12, интенсивность нагрузки и размещение: равномерное, кластерное, коридорное и регулярное. Число повторений в этих сериях определяется отдельно в настройках запуска. При изменении SF генератор размещения использует соответствующую расчётную дальность, поэтому такую серию следует интерпретировать как сравнение нормированных развёртываний, а не как изменение SF на одной фиксированной карте.

Результаты должны представляться вместе с конфигурацией, горизонтом моделирования и числом повторений. Сравнение PDR сопровождается задержкой, служебным эфиром и числом передач на доставку. Это позволяет отличать уменьшение накладных расходов от перераспределения затрат в повторные передачи данных.

\subsection{Ограничения и последующая проверка}

Ранжирование относится к заданным моделям, интервалам служебных сообщений и сценариям. Оно не переносится автоматически на все реализации AODV, OLSR или прошивки Meshtastic. Одноканальная модель с одинаковыми параметрами устройств не описывает всю неоднородность городской сети. Измерения на компьютере также не устанавливают вычислительную стоимость на встроенной платформе.

Следующий этап исследования посвящён физическому уровню действующей сети Meshtastic. Предполагается сопоставить модель с наблюдаемыми характеристиками каналов и затем проверить прогноз собственными узлами известной конфигурации. Данные шлюзов содержат только наблюдаемые ими пакеты, поэтому частоту приёма через шлюз нельзя напрямую приравнять к вероятности прямой доставки между всеми узлами. Параметры московской сети и её радиопресет должны задаваться отдельно от базового сценария таблицы~\ref{tab:params}. Завершённая натурная проверка в данной статье не заявляется.
